%!TEX program = uplatex
\RequirePackage{plautopatch}
\documentclass[uplatex,dvipdfmx]{jlreq}
\usepackage{amsmath,amssymb}

\begin{document}
\thispagestyle{empty}

\footnotetext{斎藤毅氏による英文の原稿を、止むを得ず三松・高倉が訳しました。}

\noindent
{\large\bfseries Mordell--Weil 群、代数的サイクル、モチーフ{\ldots}}
\hfill
{\large\bfseries 斎藤 毅（東大・数理）}

\bigskip

Fermat 予想の解決以後、楕円曲線の数論についての最も魅力的な未解決問題の一つは、
Birch--Swinnerton-Dyer の予想である。
それは、$\mathbb{Q}$ 上定義された楕円曲線 $E$ の Mordell--Weil ランクが、
$E$ のゼータ関数 $Z(E,s)$ の $s=1$ における零点の位数に等しいことを予言している。
すなわち、$\bmod\, p$ での $E$ の整数解の個数を Euler 積や解析接続のような
非常に超越的な方法を用いて集めることによって、
Mordell--Weil ランクという微妙な数論的不変量がとらえられる、という主張である。

\medskip

この予想は類数公式の一つの現代的な姿であり、
それはまた、現在発展中のゼータ関数の特殊値の理論を
大きく一般化するための出発点でもある。

\medskip

数体と関数体の密接な類似性により、
Tate は代数的サイクルとゼータ関数の間の関係に関する類似の予想に導かれた。
当初有限体上の代数曲面に対するものであったこの予想は、
現在 Tate 予想として広く知られている。
これは種々のコホモロジー理論（エタール、ホッジ、クリスタル{\ldots}）を備えた枠組みに
非常によく合っており、
その点をより詳しく研究するならば、我々はモチーフの世界へと導かれるであろう。

\medskip

この講演では、包括的な描像を専門以外の方にもわかりやすく伝えることを目標にする。

\end{document}